\documentclass[11pt,a4paper]{article}
\usepackage[utf8]{inputenc}
\usepackage{amsmath,amssymb,amsthm}
\usepackage{geometry}
\geometry{margin=2.5cm}
\usepackage{listings}
\usepackage{xcolor}
\usepackage{hyperref}

\lstset{
  language=C++,
  basicstyle=\ttfamily\small,
  keywordstyle=\color{blue}\bfseries,
  commentstyle=\color{green!60!black},
  stringstyle=\color{red!70!black},
  numbers=left,
  numberstyle=\tiny\color{gray},
  numbersep=8pt,
  frame=single,
  breaklines=true,
  tabsize=4,
  showstringspaces=false
}

\title{IOI 1996 -- Sorting a Three-Valued Sequence}
\author{Solution}
\date{}

\begin{document}
\maketitle

\section{Problem Statement}

Given a sequence of $n$ integers, each being 1, 2, or 3, sort the sequence in
non-decreasing order using the minimum number of swaps. Each swap exchanges exactly
two elements at arbitrary positions.

\medskip
\noindent\textbf{Constraints:} $1 \le n \le 1000$.

\section{Solution Approach}

\subsection{Zones and Misplacement Counts}

In the sorted sequence, we know exactly which positions hold each value:
\begin{itemize}
  \item \textbf{Zone 1} (positions $1 \ldots n_1$): should contain all 1s.
  \item \textbf{Zone 2} (positions $n_1{+}1 \ldots n_1{+}n_2$): should contain all 2s.
  \item \textbf{Zone 3} (positions $n_1{+}n_2{+}1 \ldots n$): should contain all 3s.
\end{itemize}
Here $n_v = |\{i : a_i = v\}|$ for $v \in \{1,2,3\}$.

Define $c_{ij}$ as the number of positions in zone $i$ that currently hold value $j$.
For example, $c_{12}$ counts the 2s sitting in positions meant for 1s.

\subsection{Two-Cycles and Three-Cycles}

A \textbf{two-cycle} is a pair of misplaced elements that can fix each other with one swap.
For instance, a 2 in zone 1 paired with a 1 in zone 2 can be swapped, fixing both.
The number of such swaps between zones $i$ and $j$ is $\min(c_{ij}, c_{ji})$.

After exhausting all two-cycles, the remaining misplaced elements form
\textbf{three-cycles}, each involving one element from each zone (e.g., a 2 in zone~1,
a 3 in zone~2, and a 1 in zone~3). Each three-cycle requires 2 swaps.

\subsection{Formula}

After performing two-cycle swaps, the residual counts satisfy
$c_{12}' = c_{23}' = c_{31}'$ (one rotation direction) and
$c_{13}' = c_{32}' = c_{21}'$ (the other direction), since the surplus must balance
across all three zones. The total number of swaps is:
\[
  \text{swaps} = \underbrace{\min(c_{12}, c_{21}) + \min(c_{13}, c_{31}) + \min(c_{23}, c_{32})}_{\text{two-cycles}}
                 + \underbrace{2 \cdot (c_{12}' + c_{13}')}_{\text{three-cycles}}.
\]

\begin{proof}[Proof of optimality]
Each swap fixes at most 2 misplaced elements. Two-cycles achieve this maximum.
Three-cycles require 2 swaps for 3 misplacements (the best possible since 1 swap
can fix at most 2). Thus no strategy can use fewer swaps.
\end{proof}

\section{C++ Solution}

\begin{lstlisting}
#include <cstdio>
#include <algorithm>
using namespace std;

int main() {
    int n;
    scanf("%d", &n);

    int a[1001];
    int cnt[4] = {}; // cnt[v] = count of value v

    for (int i = 0; i < n; i++) {
        scanf("%d", &a[i]);
        cnt[a[i]]++;
    }

    // Zone boundaries (0-indexed)
    // Zone 1: [0, cnt[1])
    // Zone 2: [cnt[1], cnt[1]+cnt[2])
    // Zone 3: [cnt[1]+cnt[2], n)

    // c[i][j] = count of value j in zone i
    int c[4][4] = {};
    for (int i = 0; i < n; i++) {
        int zone;
        if (i < cnt[1]) zone = 1;
        else if (i < cnt[1] + cnt[2]) zone = 2;
        else zone = 3;
        c[zone][a[i]]++;
    }

    // Two-cycle swaps
    int swaps = 0;
    int two12 = min(c[1][2], c[2][1]);
    int two13 = min(c[1][3], c[3][1]);
    int two23 = min(c[2][3], c[3][2]);
    swaps += two12 + two13 + two23;

    // Subtract two-cycle contributions
    c[1][2] -= two12; c[2][1] -= two12;
    c[1][3] -= two13; c[3][1] -= two13;
    c[2][3] -= two23; c[3][2] -= two23;

    // Remaining misplacements form three-cycles, each needing 2 swaps
    swaps += 2 * (c[1][2] + c[1][3]);

    printf("%d\n", swaps);
    return 0;
}
\end{lstlisting}

\section{Complexity Analysis}

\begin{itemize}
  \item \textbf{Time complexity:} $O(n)$. One pass to count values and compute zone
        membership; the rest is $O(1)$ arithmetic.
  \item \textbf{Space complexity:} $O(n)$ for the input array. The counting arrays
        use $O(1)$ additional space.
\end{itemize}

\end{document}
